\subsection{Problem statement}
At time $t$ the network is a graph $G_t=(V,E_t)$ whose vertices $V$ are the
UAVs. Each edge $e=(a,b)$ carries a measured quality $d_e(t)\in[0,1]$, the probe
reception ratio within a one-second window. A quality estimator produces the
decision variable $x_e(t)$ in one of two families:
\begin{align}
\text{measure and smooth:}\quad & x_e(t)=\alpha\, d_e(t)+(1-\alpha)\,x_e(t-1),\label{eq:ewma}\\
\text{learned forecast:}\quad & x_e(t)=\hat g_e(t)\approx
\frac{1}{H}\sum_{\tau=t+1}^{t+H} d_e(\tau),\qquad H=5\ \text{s}.\label{eq:fore}
\end{align}
Equation~\eqref{eq:ewma} is a one-parameter EWMA; $\alpha=0.5$ is the
configuration inherited from earlier studies and serves as the v8 baseline.
Equation~\eqref{eq:fore} is the GNN regressor, trained by supervised learning on
historical data and using only causal observations up to time $t$.

Both actions consume the same $x$. The routing action selects the path
\begin{equation}
\pi_t=\arg\min_{P\in\mathcal{P}(s,d)}\ \sum_{e\in P}\frac{1}{\max(x_e(t),\varepsilon)},
\label{eq:etx}
\end{equation}
that is, an ETX-style cost~\cite{decouto} evaluated on $x$ rather than on raw
measurements, under the route-pinning rule with margin and dwell described in
Section~\ref{sub:v8}. The placement action uses the best bottleneck path quality
\begin{equation}
Q(t)=\max_{P\in\mathcal{P}(s,d)}\ \min_{e\in P} x_e(t),
\label{eq:bottleneck}
\end{equation}
and triggers a fixed relay motion rule when $Q(t)$ falls below a threshold.

The central distinction of this paper lies between two measures taken on the
same $x$. \emph{Forecast} quality is measured by the mean squared error
$\mathrm{MSE}=\mathbb{E}\big[(x_e(t)-\frac{1}{H}\sum_\tau d_e(\tau))^2\big]$,
computed offline over traces. \emph{Service} quality is measured by PDR and
disconnected seconds, computed in closed loop while $x$ actually
drives~\eqref{eq:etx} or~\eqref{eq:bottleneck}. The question of the paper is
whether these two measures rank estimators in the same order.

\subsection{Environment and experimental design}
The experiments use ns-3.45~\cite{ns3}, Wi-Fi 802.11g and OLSR~\cite{rfc3626} on
chains of five or seven UAVs. Source and destination sit at the two ends of the
chain and the remaining nodes act as relays. Each episode lasts 240~s with a
256~kbps UDP main flow and two 128~kbps background flows. Two contexts are used,
static and gradual: the endpoints either stay still or separate from second 30.
The binary and runtime are inherited from the previous studies, including the
OLSR-ETX baseline implemented by the group and an ARP DeadTimeout of 1~s. Every
arm uses the same binary, channel, traffic and block seeds.

The simulation area is $1800\times1200$~m with altitudes of 50--150~m. The
initial spacing between consecutive nodes is nominally 24~m with up to 0.72~m of
jitter per coordinate. The channel combines log-distance path loss (exponent
2.2, 40~dB reference loss at 1~m) with Nakagami fading ($m=3$, $1.5$ and $1$ for
distances below 96~m, 96--180~m and above 180~m respectively). Transmit power is
20~dBm and the PHY rate is fixed at 6~Mbps. The nominal 240~m range is used to
set receiver sensitivity and is not a distance-based link cut-off. Warmup lasts
10~s; probes are broadcast every 0.05~s and aggregated over a one-second window;
the main flow uses 512~B packets, 15{,}000 packets per episode.

The protocol was written before the new simulations. Each phase packages the
source code, model, seed plan and runtime identity under SHA-256; the test phase
inherits exactly the sources and model of validation. Each cell runs once, and a
four-worker queue holds a server-wide lock. Validation comprises 80 cells (two
static and eight gradual per configuration, four arms); the test comprises 192
cells (four static and 20 gradual per configuration, four arms). No model was
reselected on the basis of validation or test service results.

\subsection{Link forecasting and the control model}
The model forecasts the mean probe reception ratio over the next five seconds,
not a link survival probability. The development data consist of 48 HOLD traces
from the v7 test phase for fitting and 20 HOLD validation traces from v7 for
epoch selection. These already-observed data are treated wholly as development
data for the new studies; v8 and v9 use independent audited seeds. Each episode
yields 236 valid forecast origins, giving 11{,}328 fitting samples and 4{,}720
epoch-selection samples; correlated frames are not treated as independent
replicates when testing control effectiveness.

The inputs comprise only current and past observations. The eight node features
are coordinates relative to the formation centroid, velocity, neighbour degree
and mean reception ratio. The seven edge features are distance, RSSI, current
reception ratio, an EWMA with coefficient 0.5, the reception-ratio difference,
the separation velocity and the magnitude of relative velocity. Coordinates and
distances are divided by 240~m, velocities by 10~m/s, and RSSI is transformed as
$(\mathrm{RSSI}+90)/40$ with missing values encoded as $-100$~dBm. An edge
participates in message passing only if its current reception ratio is at least
0.1. No future mission trajectory is read.

Edge-SAGE comprises two mean-aggregation layers with edge features, 64
dimensions and ReLU, followed by a 64-dimensional node-pair decoder with a
sigmoid and symmetrised logits. The GNN has 19{,}137 parameters. The MLP drops
the two learned message-passing branches, leaving 13{,}505 parameters, and still
retains the node-degree feature. The ablation therefore does not separate the
effect of message passing from the difference in model capacity. Both use MSE
loss, Adam with learning rate 0.001, weight decay $10^{-5}$, batch size 128,
gradient clipping at norm 5, at most 60 epochs, patience 10 and seed 20261123.
The GNN and MLP selected epochs 58 and 59. The largest discrepancy between the
NumPy and PyTorch inference paths over 40 check snapshots is
$1.19\times10^{-7}$. This is supervised learning; neither study uses
reinforcement learning.

\subsection{Routing study v8}\label{sub:v8}
The four arms are current quality (\code{lq-route}), the GNN forecast
(\code{gnn-route}), the MLP forecast (\code{mlp-route}) and OLSR-ETX-fast. The
three route-pinning arms share the same ETX rule with margin 1.0, a three-step
warmup and a five-second dwell; the broken-link check still uses a current
reception ratio below 0.1. ETX-fast uses the HELLO interval of 0.5~s and TC
interval of 1.25~s selected in v7. Relays stay still, and each gradual endpoint
has a nominal speed of 1~m/s multiplied by a factor in $[0.85, 1.15]$ drawn per
layout seed. The source moves along $+y$ and the destination along $-y$; the
same seed yields the same endpoint motion across arms. Validation checks
technical criteria only: cell count, route correctness, simulator flags,
forecasting, counters and background flows, no startup blackout and no relay
motion. The go decision does not depend on any PDR difference.

The three tests in fixed order are the PDR difference of GNN minus current, GNN
minus MLP, and then GNN minus ETX-fast. A step succeeds when the lower bound of
the 95\% confidence interval is positive; an affirmative interpretation is
permitted only when every preceding step has succeeded. Test MSE is measured
separately for the GNN, the MLP and the EWMA on the same \code{lq-route} traces,
kept apart from closed-loop PDR.

\subsection{Placement study v9}
Study v9 keeps the v8 weights unchanged and routes on current quality in every
arm. The arms are stationary relays, reactive motion triggering, GNN triggering
and MLP triggering. The best path quality is $Q=\max_P\min_{e\in P}q_e$. After
three frames, the reactive arm triggers if the current EWMA has $Q<0.70$ in two
consecutive frames; the GNN and MLP arms add the analogous condition on the
forecast. A forecast edge is used only if its current measurement is at least
0.1.

Relays are ordered by their projection onto the line joining the current
endpoints and then directed towards points that divide that segment equally. The
limits are 8~m/s, 400~m per relay, a 1~m deadband, at most ten steps per burst
and a five-step rest. Motion stops early after three frames in which the current
quality, and the forecast where applicable, are at least 0.80. An in-step
minimum separation check must reach 5~m; otherwise the relay motion command is
set to zero. This rule is identical in every arm with motion, only the trigger
differs.

The nominal speed of each gradual endpoint rises to 2~m/s, retaining the
per-seed random factor of v8, in order to create a disconnection challenge; this
differs from the 1~m/s fitting distribution. Validation requires the technical
and safety criteria to pass and at least 8 of 16 stationary gradual baseline
episodes to lose connectivity for at least five seconds. If that fails, the test
does not run and the speed is not retuned within the same study.

The primary endpoint is the disconnected seconds of the reactive arm minus those
of the GNN arm; the subsequent steps consider the PDR of GNN minus reactive, and
then the disconnected seconds of MLP minus GNN. Disconnection means that source
and destination lie in different connected components of the probe graph
thresholded at 0.5, counted over frames 1--240. The seconds during which OLSR
has no route, and PDR, are separate quantities. This is a measured
graph-threshold partition, not a proof that no radio packet whatsoever could
cross between components. Travelled distance, longest outage, number of motion
bursts and background-flow PDR are descriptive.

\subsection{Statistical analysis and cost}
Each comparison keeps all four arms within the same seed block; the bootstrap
draws 10{,}000 resamples over the 40 gradual test blocks, 20 per configuration,
balanced between n5 and n7, with RNG seed 20261120. The
bootstrap~\cite{efron1993} is applied at the scenario level so that correlations
between arms are preserved. Negative results and gated steps are still
reported. A confidence interval containing zero does not establish equivalence.
Static results are descriptive only, with eight blocks.

The collector measures whole-decision latency on an eight-vCPU KVM virtual
machine (Intel Xeon Silver 4110 at 2.10~GHz) running four workers, each with
compute libraries restricted to a single thread. Probes are emitted in every
arm; centralized telemetry and commands are not yet transmitted as
channel-contending packets. Computation time is not converted into end-to-end
control delay, bytes are not converted into airtime, and travelled distance is
not converted into energy.

Signalling overhead was measured in study v7 with the same binary and runtime,
so the following figures apply to the v8 configuration: about 0.29~kB/s for
plain OLSR, 0.44~kB/s for OLSR-ETX, 1.78~kB/s for OLSR-ETX with 0.5~s HELLO, and
about 5.55~kB/s for route pinning, comprising 0.29~kB/s of fallback OLSR,
5.02~kB/s of probes and an estimated 0.25~kB/s of reports to the controller.
These are IP-layer bytes plus the estimated reporting component and are not
converted into a percentage of channel capacity. Probes are emitted in every
arm, so probe airtime is common across arms, but it is charged only to the
methods that consume probes. The centralized reporting component does not
contend for the channel in simulation. Every comparison between route pinning
and OLSR-ETX in this paper must therefore be read as an upper bound for a
partially idealized centralized controller.
